\section{模型的评价与推广}

\subsection{模型的优点}

\begin{enumerate}[label=\textbf{优点 \arabic*}]
	\item \textbf{物理口径唯一，四问横向可比。}附录给定的时间、能耗、充电与链路公式只在公共物理层实现一次（航段几何、载荷—航程、逐段能耗、两阶段充电、链路预算及三态判定），问题一至问题四全部复用，杜绝了“各问各写一套”造成的口径分叉。这是四问结果能够互相印证、层层继承的前提。

	\item \textbf{数值反解严谨且可验证。}最大安全载荷因指数 $3/2$ 无初等反函数，本文用单调性 $+$ Brent 法数值反解并\textbf{强制回代验证}约束紧性，实测相对误差 $5.6\times10^{-15}$，达到机器精度。工程上“以近似公式代替反解”是常见失分点，本文避免了这一陷阱。

	\item \textbf{结论可证最优而非仅“较优”。}问题一通过解析下界证明了 18 架次\textbf{在架次数维度上最优}（15 个服务区逐区等于下界）；问题四的合法组数判据是\textbf{精确的图论结论}，不含启发式近似。这两处使全文有两条“硬结论”支撑。

	\item \textbf{结果可独立复算。}独立校验器\textbf{不导入任何求解器模块}，从附件与 DEM 重算全部物理量；本文另用 240 条航段缓存对问题二全部 25 个架次\textbf{逐段重算}，与求解器上报值一致（总能耗均为 \SI{77.3016}{kWh}），资源占用零冲突、结束时刻自洽偏差最大 \SI{0.078}{s}。\textbf{两套独立计算}显著提升了结果可信度。

	\item \textbf{勇于纠正否定性结论。}本文两次纠正了会改变结论方向的建模缺陷（\S\ref{sec:q2traps}、\S\ref{sec:q3bug}）：一次推翻了“首批时限物理不可行”，一次推翻了“通信覆盖 100\%”。后者尤其值得强调——\textbf{“校验器报 0 违规”不等于“方案没有问题”}，必须先确认校验器确实执行了该项检查。这种自我复核过程本身就是建模严谨性的体现。

	\item \textbf{关键结论具备工程指导性。}如“\textbf{结构载重而非能量才是主导约束}”、“\textbf{中继接入段最弱故应贴近作业空域而非网关}”、“\textbf{异构机队必须真正混编}（同构 C 型方案准时率仅 47.5\%，混编后达 100.0\%）”、“\textbf{分区不省资源但可削峰}（资源总量 33/34/38，最大组工作量 13.95$\to$13.00\,h）”、“\textbf{压架次数不能靠调权重}”、“\textbf{覆盖率必须用时间轴口径}”等，均给出了可操作的量化依据。
\end{enumerate}

\subsection{模型的缺点与不足}

\begin{enumerate}[label=\textbf{不足 \arabic*}]
	\item \textbf{问题二的最优性未证明。}25 架次距理论下界 24 仅高 1 个，但本文只称其为可行解。若要以精确方法证明最优，需把组批做成\textbf{集合划分整数规划}（每箱恰好覆盖一次，目标函数为架次数），并用\textbf{状态压缩动态规划}在小规模服务区上交叉验证；本文采用启发式是出于求解规模的考虑，属于\textbf{精度与效率的权衡}。

	\item \textbf{中继选址采用充分条件而非最优控制。}式~\eqref{eq:suffcond} 是\textbf{充分条件}——只要找到点就必然可行，但可能遗漏“需要多次切换中继位置”的更优方案。若升级为最优控制（每时刻可切换中继与悬停点的动态规划），有望降低中继架次数。

	\item \textbf{问题三的覆盖率受题目资源约束封顶。}55.0\% 的时间轴覆盖率并非算法能力不足，而是 \S\ref{sec:q3gap} 所述的资源缺口所致。本文如实报告而非以几何可达覆盖率替代，但这也意味着该问的结论\textbf{依赖题目给定的中继数量}。

	\item \textbf{未考虑动态因素。}模型假定任务期间天气、风场、临时禁飞区与通信干扰恒定；实际救援中这些因素可能显著改变可行域。
\end{enumerate}

\subsection{推广}

本文的模型框架可直接推广到更一般的\textbf{多基地异构机队应急投送}场景：

\textbf{（1）方法层面。}“\textbf{载荷—航程非线性 $+$ 二维装箱 $+$ 解析下界}”这一组合适用于任何“载具能力随载荷衰减”的投送问题；把“\textbf{任务分区建模为图连通分量}”的思路适用于任何“同一执行单元涉及的对象必须同组”的分区问题，例如共享车辆的多站点配送分组。\textbf{“真实调度器复核每一次合并”}这一做法尤其值得推广——它是避免“目标函数与上报口径分叉”的通用手段。

\textbf{（2）场景层面。}把单基地 $O01$ 扩展为多基地时，只需在载荷层加入“最近基地选择”，其余分层结构不变；把中继无人机扩展为可移动中继，则式~\eqref{eq:suffcond} 的充分条件需升级为时空联合优化，但“\textbf{覆盖率必须区分几何可达与时间轴在站}”这一判据仍然成立。

\textbf{（3）结论层面。}本文得到的“结构载重主导”“接入段最弱”“机队必须混编”“分区削峰不省资源”“中继资源是短板”等结论，对同类山区应急救援场景具有参考价值；其中\textbf{“资源缺口应如实报告而非用宽松口径掩盖”}是应急优化建模中最重要的方法论结论。
